% !PRINTED 611
\weberpage{611}
We now form the field:

\weq{16}{\mathfrak{K}'' = \Omega(\xi, \xi') = \Omega(\xi''),}

by setting

\weq{17}{\xi'' = \alpha\xi + \alpha'\xi'}

and we denote by $\mathfrak{P}''$ a prime ideal of the first degree of this field
that divides some $\mathfrak{p}_1'$.

Let $\mathfrak{P}, \mathfrak{P}'$ be, as above, the prime ideals in
$\mathfrak{K}, \mathfrak{K}'$ divisible by $\mathfrak{P}''$. Then by (15) and (17)
\[
\xi''^{p} \equiv \xi'' \quad (\mathrm{mod}\ \mathfrak{P}''),
\]
and consequently $\mathfrak{K}''$ is a class field of $A'$.

Accordingly $\mathfrak{K}''$ is identical with $\mathfrak{K}'$; $\xi'$ is a primitive number
of $\mathfrak{K}''$, and consequently every number in $\mathfrak{K}''$, and hence $\xi$ as well, can be
expressed rationally by $\xi'$.

From this it follows:

11. \gesperrt{If a number group $A$ contains another number
 group $A'$ as a divisor of finite index, and if the
 class fields $\mathfrak{K}(A)$, $\mathfrak{K}(A')$ exist, then $\mathfrak{K}(A)$
is contained in $\mathfrak{K}(A')$.}

\begin{center}
\textbf{§ 165. Prime Ideals in the Classes.}
\end{center}

We now return to the functions $Q$ which we represented in
§ 163, (16) by infinite series and in § 163, (18) by
infinite products:

\weq{1}{Q_{\varkappa} = \Pi\, \frac{1}{1 - \chi_{\varkappa}(\mathfrak{p})\,N(\mathfrak{p})^{-s}}}

in which $\chi_{\varkappa}$ denotes one of the characters of the group
and $\mathfrak{p}$ runs through all prime ideals in $\overline{O}$. We develop the logarithm
 of this product in the following way:

\weq{2}{\log Q_{\varkappa} = \sum \frac{\chi_{\varkappa}(\mathfrak{p})}{N(\mathfrak{p})^{s}} + \frac{1}{2}\sum \frac{\chi_{\varkappa}(\mathfrak{p}^2)}{N(\mathfrak{p})^{2s}} + \frac{1}{3}\sum \frac{\chi_{\varkappa}(\mathfrak{p}^3)}{N(\mathfrak{p})^{3s}} + \cdots}

We do not need to know the multiple of $2\pi i$ that is not more closely defined
in the $\log Q$, since it could change with $s$ only discontinuously; and since
continuous functions of $s$ stand on the right and on the left, as long
as $s > 1$, this multiple does not change with $s$ at
 all.

Let $\overline{A}_i$ be one of the classes of the ideal group $\overline{A}$ and $\overline{A}_i'$ the
opposite class. Then by (12), (13), (14), § 163:
\[
\overset{k}{\sum} \chi_{\varkappa}(\overline{A}_i')\,\chi_{\varkappa}(\mathfrak{a}) = h \ \text{or} = 0,
\]
%CONT>>
